\documentclass[11pt]{article}
\usepackage[T1]{fontenc}
\usepackage{lmodern}
\usepackage{microtype}
\usepackage[margin=0.82in]{geometry}
\usepackage{amsmath,amssymb,amsthm}
\usepackage{xcolor}
\usepackage{hyperref}
\usepackage{enumitem}
\usepackage{parskip}

\definecolor{navy}{HTML}{183153}
\definecolor{teal}{HTML}{167D7F}
\definecolor{pale}{HTML}{EEF6F6}
\definecolor{ink}{HTML}{20242A}
\hypersetup{colorlinks=true,linkcolor=teal,urlcolor=teal,citecolor=teal}
\color{ink}
\setlist{nosep,leftmargin=1.5em}
\newtheorem{theorem}{Theorem}

\newcommand{\R}{\mathbb R}
\newcommand{\dd}{\,d}
\newcommand{\Ltwo}[1]{L^2\!\left(G,#1\,dg\right)}

\begin{document}

{\color{navy}\LARGE\bfseries A Harnack answer to Mori's Problem 49}\par
\vspace{0.25em}
{\large Later-time heat-kernel $L^2$ spaces are dense}\par
\vspace{0.55em}
\noindent\colorbox{pale}{\parbox{0.965\linewidth}{
\textbf{Status.} Literature-implied answer to the full problem. The decisive
input is Li--Xu's 2009 parabolic Harnack theorem; the implication below does
not appear to be stated explicitly in the inspected literature.
}}

\section*{Source and exact question}

Shota Mori, \emph{The heat kernel on $SL(2,\R)$}, arXiv:1909.03670v3
\cite{Mori}, Section 8, Problem 49, PDF page 25:

\begin{quote}
Let $t,\varepsilon>0$. Prove that the space
\[
L^2\!\left(SL(2,\R),\rho_{SL(2,\R)}(t+\varepsilon,g)\,dg\right)
\]
is a dense subspace of
\[
L^2\!\left(SL(2,\R),\rho_{SL(2,\R)}(t,g)\,dg\right).
\]
\end{quote}

\section*{Identification}

Mori's metric is left invariant. Hence it is complete, and its Ricci tensor
has a global lower bound $\operatorname{Ric}\ge -k g$ for some $k\ge0$:
left translations are isometries, so the Ricci eigenvalues are the same at
every point. Li and Xu's parabolic Harnack inequality
\cite[Theorem~1.4]{LiXu} applies to every positive heat solution on such a
complete noncompact manifold.

At the same spatial point $x_1=x_2=g$ and times $t_1=t<t_2=s$, their theorem
gives
\begin{equation}\label{eq:ratio}
 \rho_t(g)\le C_{t,s}\rho_s(g)\qquad(g\in G),
\end{equation}
where, if $n=\dim G$,
\[
C_{t,s}=\left(\frac{s}{t}\right)^{n/2}
\left(\frac{1+2ks/3}{1+2kt/3}\right)^{-n/8}
\exp\!\left(\frac{nk}{4}(s-t)\right)<\infty.
\]
The distance term in the Harnack inequality vanishes, so the constant is
uniform in $g$.

\begin{theorem}
Let $G$ be a connected noncompact Lie group with a left-invariant Riemannian
metric, and let $\rho_t$ be its positive heat kernel for
$\partial_tu=\Delta u$. For $0<t<s$, the natural inclusion
\[
\Ltwo{\rho_s}\hookrightarrow\Ltwo{\rho_t}
\]
is continuous and has dense range. In particular, Mori's Problem 49 has an
affirmative answer by taking $G=SL(2,\R)$ and $s=t+\varepsilon$.
\end{theorem}

\begin{proof}
The Harnack comparison \eqref{eq:ratio} gives, for every $f\in\Ltwo{\rho_s}$,
\[
 \|f\|_{L^2(\rho_tdg)}^2
 =\int_G|f(g)|^2\rho_t(g)\,dg
 \le C_{t,s}\int_G|f(g)|^2\rho_s(g)\,dg.
\]
Thus the natural inclusion is well-defined and continuous. Both heat kernels
are positive and continuous. Consequently $C_c(G)\subset\Ltwo{\rho_s}$,
while $C_c(G)$ is dense in $\Ltwo{\rho_t}$ for the Radon measure
$\rho_t(g)\,dg$. Hence the range of the inclusion is dense.
\end{proof}

\section*{Provenance, scope, and verification}

This is an agent-identified implication, not an original theorem claimed by
this run. Li--Xu predates Mori's question by ten years and therefore does not
identify Problem 49. The proof resolves Problem 49 completely and in fact
works on every connected noncompact Lie group with a left-invariant metric.
It makes no claim about Mori's Problems 46--48. For compact connected Lie
groups, the same conclusion follows directly because the ratio of two
positive continuous heat kernels is bounded on the compact group.

The hypothesis audit checks completeness, the Ricci lower bound, positivity,
and matching heat-equation normalization. A bounded search of the run indexes,
the exact wording of Problem 49, arXiv:1909.03670, and author/title variants
found no paper explicitly stating this answer. The official source record is
still v3 (2 October 2019). Human review should focus on the application of
Li--Xu Theorem 1.4 to Mori's normalization; both use
$\partial_tu=\Delta u$ in the relevant statements.

\begin{thebibliography}{9}
\bibitem{Mori}
S. Mori, \emph{The heat kernel on $SL(2,\mathbb R)$}, arXiv:1909.03670v3
(2019), Section 8, Problem 49.

\bibitem{LiXu}
J. Li and X. Xu, \emph{Differential Harnack inequalities on Riemannian
manifolds I: linear heat equation}, arXiv:0901.3849 (2009), Theorem 1.4.
\end{thebibliography}

\end{document}
